
\noindent
Zur Bestimmung der maximalen Leistung wird die Abhängigkeit der Lastleistung vom Lastwiderstand untersucht. 
Ausgangspunkt ist die Gleichung~\ref{eq:Leistung_Last_Spannungsquelle}.

\noindent
Um das Maximum der Leistung zu bestimmen, wird die erste Ableitung der Leistung nach dem Lastwiderstand $R_L$ gebildet:
\begin{equation}
\label{eq:dPL_dRL}
P_L'(R_L)
= U_0^2 \frac{(R_I + R_L)^2 - 2R_L(R_I + R_L)}{(R_I + R_L)^4}
\end{equation}

\noindent
Der Zähler kann vereinfacht werden zu:
\begin{equation}
\label{eq:dPL_dRL_simpl}
P_L'(R_L)
= U_0^2 \frac{R_I - R_L}{(R_I + R_L)^3}
\end{equation}

\noindent
Das Maximum der Leistung liegt an der Stelle, an der die Ableitung gleich null ist. Da der Nenner der Ableitung stets positiv ist, bestimmt allein der Zähler die Nullstelle:
\begin{equation}
\label{eq:nullstelle_zaehler}
R_I - R_L = 0
\end{equation}

\noindent
Daraus folgt unmittelbar:
\begin{equation}
\label{eq:RL_eq_RI}
R_L = R_I
\end{equation}

\noindent
Zur Absicherung, dass es sich hierbei tatsächlich um ein Maximum handelt, kann das Vorzeichen der Ableitung betrachtet werden. Für $R_L < R_I$ ist die Ableitung positiv, während sie für $R_L > R_I$ negativ ist. Die Leistung besitzt daher an dieser Stelle ein eindeutiges Maximum.

\noindent
Zur Klassifikation der Extremstelle wird zusätzlich die zweite Ableitung der Leistung $P_L(R_L)$ nach dem Lastwiderstand betrachtet:
\begin{equation}
\label{eq:d2PL_dRL2}
P_L''(R_L)
= -\frac{2U_0^2}{(R_I + R_L)^3}
\end{equation}

\noindent
Da der Nenner für $R_I + R_L > 0$ positiv ist, ist die zweite Ableitung negativ. Insbesondere gilt an der Stelle $R_L = R_I$:
\begin{equation}
\label{eq:d2PL_negative_at_max}
P_L''(R_L) < 0
\end{equation}

\noindent
Damit handelt es sich bei der zuvor bestimmten Extremstelle eindeutig um ein Maximum der Leistung.

\noindent
Durch Einsetzen der Bedingung $R_L = R_I$ in die Gleichung~\ref{eq:Leistung_Last_Spannungsquelle} lässt sich eine allgemeine Formel zur Berechnung der maximalen Leistung herleiten:
\begin{equation}
\label{eq:P0}
P_0 = \frac{U_0^2}{4R_I}
\end{equation}

\noindent
Setzt man die gegebenen Werte $U_0 = 15\,\mathrm{V}$ und $R_I = 5\,\Omega$ ein, ergibt sich für die maximale Leistung der Spannungsquelle:
\begin{equation*}
\label{eq:P0_numeric}
P_0 = \frac{(15\,\mathrm{V})^2}{4 \cdot 5\,\Omega} = 11.25\,\mathrm{W}
\end{equation*}
